% !TeX program = xelatex
\documentclass[11pt,a4paper]{ctexart}

\usepackage[a4paper,margin=2.2cm]{geometry}
\usepackage{amsmath,amssymb,bm}
\usepackage{booktabs}
\usepackage{array}
\usepackage{hyperref}
\usepackage{xcolor}
\usepackage{graphicx}

\hypersetup{
  colorlinks=true,
  linkcolor=blue,
  urlcolor=blue,
  citecolor=blue
}

\newcommand{\Tr}{\operatorname{Tr}}

\title{\textbf{氘核自旋密度矩阵与 $n$--$p$ 空间分布}\\
\large SAMURAI 竖直纯张量极化束流}
\author{}
\date{}

\begin{document}
\maketitle

\begin{abstract}
氘核的 $S$--$D$ 波耦合使自旋极化与质子--中子相对坐标的角分布相关。
对于绕竖直 $y$ 轴轴对称的束流，空间密度的四极项正比于
$p_{yy}[\sqrt2\,uw-w^2/2]P_2(\cos\theta_y)$，其中 $u,w$ 为径向波函数。
在该径向系数为正的区域，$p_{yy}=+1$ 使沿 $y$ 轴的概率增强，
$p_{yy}=-2$ 使垂直于 $y$ 的平面内概率增强。
两种模式具有相同的径向距离分布。
这一关系由束流的自旋密度矩阵与氘核波函数共同确定，
可用于半经典模型的坐标抽样。
\end{abstract}

\section{氘核波函数中的自旋与轨道耦合}

氘核基态的量子数为 $J^\pi=1^+$。
在非相对论描述中，基态以 ${}^3S_1$ 成分为主，
核力的张量部分将其与 ${}^3D_1$ 成分耦合。
选定量子化轴后，总角动量投影为 $M$ 的波函数写为
\begin{equation}
\Psi_{1M}(\bm r,\sigma)
=
\frac{u(r)}{r}Y_{00}(\hat{\bm r})\chi_{1M}(\sigma)
+
\frac{w(r)}{r}
\sum_{m_L,m_S}
\langle 2m_L,1m_S|1M\rangle
Y_{2m_L}(\hat{\bm r})\chi_{1m_S}(\sigma),
\label{eq:deuteronwf}
\end{equation}
其中 $\bm r=\bm r_p-\bm r_n$ 为质子--中子相对坐标，
$\sigma$ 表示两个核子的自旋自由度，$\chi_{1m_S}$ 为 $S=1$ 自旋三重态。
$S$ 波和 $D$ 波的径向函数分别为 $u(r)$ 和 $w(r)$，满足
\[
\int_0^\infty [u^2(r)+w^2(r)]\,dr=1.
\]

在纯 $S$ 波极限 $w=0$ 下，空间部分 $Y_{00}$ 球对称且与 $M$ 无关，
因而任意磁子态占据都给出球对称的空间分布。
加入 $D$ 波后，Clebsch--Gordan 系数将 $m_L$ 与 $m_S$ 关联起来，
总波函数一般无法分解为独立的空间部分和自旋部分。
这种耦合使不同磁子态具有不同的空间角分布。

\section{磁子态占据与束流极化}

取竖直 $y$ 方向为量子化轴，采用基底
$\{|+1\rangle_y,|0\rangle_y,|-1\rangle_y\}$。
束流的 $3\times3$ 自旋密度矩阵描述这三个 $J=1$ 磁子态的占据和相干：
\begin{equation}
\rho^{(J=1)}
=
\begin{pmatrix}
\rho_{++} & \rho_{+0} & \rho_{+-}\\
\rho_{0+} & \rho_{00} & \rho_{0-}\\
\rho_{-+} & \rho_{-0} & \rho_{--}
\end{pmatrix},
\qquad
\Tr\rho^{(J=1)}=1.
\label{eq:spinrho}
\end{equation}
绕 $y$ 轴轴对称时，密度矩阵在该基底下为对角形式：
\begin{equation}
\rho^{(J=1)}=\operatorname{diag}(n_+,n_0,n_-),
\qquad n_++n_0+n_-=1.
\end{equation}
矢量极化和张量极化定义为
\begin{align}
p_y &= n_+-n_-,\\
p_{yy} &= n_++n_--2n_0,
\label{eq:polarization}
\end{align}
取值范围分别为 $-1\le p_y\le1$ 和 $-2\le p_{yy}\le1$。
纯张量极化满足 $p_y=0$、$p_{yy}\ne0$。

\section{对核子自旋取迹，得到空间密度}

将束流密度矩阵与式~\eqref{eq:deuteronwf} 的完整波函数结合，
得到包含轨道和核子自旋自由度的统计算符
\begin{equation}
\hat\rho_{\rm tot}
=
\sum_{M,M'=-1}^{1}
\rho_{MM'}^{(J=1)}|\Psi_M\rangle\langle\Psi_{M'}|.
\label{eq:fullrho}
\end{equation}
对未观测的核子自旋取迹，得到约化空间密度矩阵
\begin{equation}
\rho_{\rm sp}(\bm r,\bm r')
=
\sum_\sigma
\langle\bm r,\sigma|\hat\rho_{\rm tot}|\bm r',\sigma\rangle.
\label{eq:reduced}
\end{equation}
其对角元给出单位体积内的概率密度：
\begin{equation}
\varrho(\bm r)
=
\rho_{\rm sp}(\bm r,\bm r)
=
\sum_{M,M'}\rho_{MM'}^{(J=1)}
\sum_\sigma\Psi_M(\bm r,\sigma)\Psi_{M'}^*(\bm r,\sigma).
\label{eq:localdensity}
\end{equation}
因此，空间分布由磁子态占据、相干性和氘核内部波函数共同决定。
$\rho_{\rm sp}$ 也适用于混合束流；即使完整氘核态为纯态，
轨道与自旋的纠缠也可使取迹后的空间态成为混合态。

\section{轴对称束流的空间分布}

定义
\[
\cos\theta_y=\hat{\bm r}\cdot\hat{\bm y},
\qquad P_2(\cos\theta_y)=\frac12(3\cos^2\theta_y-1),
\]
以及径向四极系数
\begin{equation}
A(r)=\sqrt2\,u(r)w(r)-\frac12w^2(r).
\label{eq:A}
\end{equation}
将式~\eqref{eq:deuteronwf} 代入式~\eqref{eq:localdensity}，
对核子自旋求和，得到各磁子态的空间密度：
\begin{align}
\varrho_0(r,\theta_y)
&=\frac{u^2+w^2-2A(r)P_2(\cos\theta_y)}{4\pi r^2},
\label{eq:rho0}\\
\varrho_{\pm1}(r,\theta_y)
&=\frac{u^2+w^2+A(r)P_2(\cos\theta_y)}{4\pi r^2}.
\label{eq:rhopm}
\end{align}
按占据数加权，$\varrho=n_+\varrho_{+1}+n_0\varrho_0+n_-\varrho_{-1}$，
可得
\begin{equation}
\boxed{
\varrho(r,\theta_y)
=\frac{1}{4\pi r^2}
\left[u^2(r)+w^2(r)+p_{yy}A(r)P_2(\cos\theta_y)\right].
}
\label{eq:master}
\end{equation}
由于 $\varrho_{+1}=\varrho_{-1}$，加权结果只含 $n_++n_-$，
矢量极化 $p_y=n_+-n_-$ 因而从标量坐标密度中消去。
自旋相关散射仍可对矢量极化敏感。

\subsection{\texorpdfstring{$S$--$D$}{S--D} 干涉决定四极项的强度}

$A(r)$ 同时包含干涉项 $\sqrt2\,uw$ 和 $D$ 波项 $-w^2/2$。
在 $|w|\ll|u|$ 的区域，干涉项是一阶于 $w$ 的贡献，
因此较小的 $D$ 波概率也能产生明显的角分布变化。
式~\eqref{eq:deuteronwf} 中的 $S$、$D$ 成分相干叠加，
计算空间密度时需要保留这一干涉项。

\subsection{径向分布与未极化极限}

利用 $\int P_2(\cos\theta_y)\,d\Omega=0$，对式~\eqref{eq:master} 作角积分，得到
\begin{equation}
r^2\int d\Omega\,\varrho(r,\Omega)=u^2(r)+w^2(r).
\label{eq:radial}
\end{equation}
因此，对于同一组束缚态波函数 $u,w$，任意磁子态占据均有
$P(r)\,dr=[u^2(r)+w^2(r)]\,dr$。
极化通过四极项改变给定 $r$ 处的方向概率，径向距离分布保持不变。

未极化束流满足 $n_+=n_0=n_-=1/3$，故 $p_y=p_{yy}=0$。
三个磁子态的四极项相互抵消，空间密度为球对称的
$\varrho=(u^2+w^2)/(4\pi r^2)$。

\section{两种理想纯张量模式}

两种模式均满足 $p_y=0$，其占据数和空间特征如下。
表中的方向增强均以 $A(r)>0$ 为条件。
\begin{center}
\begin{tabular}{>{\centering\arraybackslash}p{2.3cm}
                >{\centering\arraybackslash}p{3.4cm}
                p{7.1cm}}
\toprule
束流状态 & $(n_+,n_0,n_-)$ & 空间分布（$A(r)>0$）\\
\midrule
$p_{yy}=+1$ & $(1/2,0,1/2)$ & 沿 $y$ 轴增强，长椭球型（prolate）\\
$p_{yy}=-2$ & $(0,1,0)$ & 在 $x$--$z$ 平面增强，扁椭球型（oblate）\\
未极化 & $(1/3,1/3,1/3)$ & 四极项抵消，球对称\\
\bottomrule
\end{tabular}
\end{center}

$p_{yy}=+1$ 对应 $M_y=+1$ 与 $M_y=-1$ 的等权非相干混合，
$p_{yy}=-2$ 对应纯态 $|J=1,M_y=0\rangle$。
这里“纯张量”限定的是极化分量；束流的量子态纯度由密度矩阵决定。

沿 $y$ 轴有 $P_2(1)=1$，在 $x$--$z$ 平面内有 $P_2(0)=-1/2$。
将这两个值代入式~\eqref{eq:master}，即可得到表中的方向增强关系。
若某个径向区域内 $A(r)<0$，相应的增强方向反转。
长椭球型和扁椭球型描述给定 $r$ 处的角分布各向异性，
两种极化模式的平均 $n$--$p$ 距离相同。

\clearpage
\section{两种模式的概率密度对比}

图~\ref{fig:rnp-density} 对比 $p_{yy}=+1$ 和 $-2$ 的相对坐标概率密度：
空间各向异性随极化改变，径向曲线则完全重合。
为直观展示式~\eqref{eq:master}，采用以下归一化示意模型：
\begin{equation}
u(r)=\frac{2\sqrt{1-P_D}}{\pi^{1/4}b^{3/2}}r e^{-r^2/(2b^2)},
\qquad
w(r)=\frac{4\sqrt{P_D}}{\sqrt{15}\,\pi^{1/4}b^{7/2}}r^3 e^{-r^2/(2b^2)},
\label{eq:plot-model}
\end{equation}
其中 $b=2\,\mathrm{fm}$，$P_D=\int w^2\,dr=0.05$。
这些参数用于示意，并非由特定核势拟合得到的氘核波函数。
令 $r_{np}=|\bm r_p-\bm r_n|=r$、$\mu=\cos\theta_y$，
对半径和方位角积分后，角度边缘概率密度为
\begin{equation}
\frac{dP}{d\mu}=\frac12[1+p_{yy}C P_2(\mu)],\qquad
C=\int_0^\infty A(r)\,dr
=\sqrt{\frac65P_D(1-P_D)}-\frac{P_D}{2}.
\end{equation}
绘图程序为 \texttt{plots/plot\_rnp\_density.py}；
在本目录用 \texttt{python3} 运行该程序，即按上述默认参数生成图像。

\begin{figure}[htbp]
\centering
\includegraphics[width=0.94\textwidth]{plots/output/rnp_density_pyy.pdf}
\caption{$n$--$p$ 概率密度示意。
（a,b）三维密度在 $z_{np}=0$ 处的截面，共用单位为 $\mathrm{fm}^{-3}$ 的色标；
这不是沿 $z$ 积分的投影。
（c）两种模式的径向密度 $P(r_{np})=u^2+w^2$ 完全相同。
（d）$\mu=\cos\theta_y$ 的归一化概率密度区分沿 $y$ 轴
（$\mu=\pm1$）与垂直平面内（$\mu=0$）的方向概率。}
\label{fig:rnp-density}
\end{figure}

\clearpage
\section{磁子态相干与方位角依赖}

非对角元 $\rho_{MM'}$（$M\ne M'$）表示磁子态之间的相干。
其中的张量分量可产生方位角 $\phi$ 依赖；
这类分布需由完整的式~\eqref{eq:localdensity} 计算，
式~\eqref{eq:master} 的单一 $p_{yy}P_2(\cos\theta_y)$ 项只适用于轴对称情形。

一般的自旋 $1$ 密度矩阵可分解为 rank-0、rank-1 和 rank-2 分量，
分别对应归一化、矢量极化和张量极化。
轴对称纯张量束流保留归一化分量及沿所选轴的张量分量。

\clearpage
\section{QMD 与半经典模型的坐标抽样}

轴对称束流可直接按式~\eqref{eq:master} 抽样相对坐标。
由于 $\varrho$ 是单位体积的密度，球坐标下的概率元为
\begin{equation}
 dP=\varrho(r,\theta_y)r^2\,dr\,d\Omega
 =\frac{u^2+w^2+p_{yy}A(r)P_2(\cos\theta_y)}{4\pi}\,dr\,d\Omega,
\label{eq:sampling}
\end{equation}
其中 $d\Omega=\sin\theta_y\,d\theta_y\,d\phi$。
可先按 $P(r)=u^2+w^2$ 抽样 $r$，再按
\[
p(\Omega\mid r)=\frac{1}{4\pi}
\left[1+\frac{p_{yy}A(r)}{u^2(r)+w^2(r)}P_2(\cos\theta_y)\right]
\]
抽样方向；轴对称性使 $\phi$ 在 $[0,2\pi)$ 上均匀分布。

若将所有相对坐标固定沿 $y$ 轴或限制在 $x$--$z$ 平面，
就会人为强化方向约束，可能放大计算中的极化效应。
上述抽样保留了量子态给出的坐标概率。
对于显式依赖核子自旋的反应模型，还需保留
式~\eqref{eq:deuteronwf} 中 $m_L$ 与 $m_S$ 的关联。

\end{document}
